\documentclass[11pt,reqno]{amsart}
\usepackage{amsmath,amssymb,graphicx,amscd}
\usepackage{mathrsfs}
\usepackage{enumerate}
% pour les S des spheres
\usepackage{mathrsfs}
\usepackage{dsfont}
\usepackage{stmaryrd}
\usepackage{wrapfig}
\usepackage{pstricks}

% ----------------------------------------------------------------
\vfuzz2pt % Don't report over-full v-boxes if over-edge is small
\hfuzz2pt % Don't report over-full h-boxes if over-edge is small

\newcommand{\prob}{{\mathbb{P}}}

\renewcommand{\S}{\mathscr{S}}  % complex unit sphere

\newcommand{\CC}{{\mathbb C}}
\newcommand{\HH}{{\mathbb H}}
\newcommand{\RR}{{\mathbb R}}
\newcommand{\QQ}{{\mathbb{Q}}}
\newcommand{\NN}{{\mathbb{N}}}
\newcommand{\ZZ}{{\mathbb{Z}}}
\renewcommand{\P}{{\mathbb{P}}}
\renewcommand{\span}{{\mathrm{span}}}
\newcommand{\sgn}{\operatorname{{sgn}}}
 \DeclareMathOperator{\Prob}{\mathbb{P}}   %probability
 \DeclareMathOperator{\E}{\mathbb{E}}      %expectation
 \DeclareMathOperator{\cov}{cov}             %logarithmic integral
\newcommand{\1}{\mathds{1}}

 \newcommand{\I}{1\!\!1}                   %Indicator function


 \newcommand{\rd}{{\mathrm{d}}}            %differential
 \newcommand{\ii}{{\mathrm{i}}}
  \newcommand{\Id}{{\mathrm{Id}}}
\newcommand{\law}{\overset{\mbox{\rm \scriptsize law}}{=}}
\newcommand{\convlaw}{\overset{\mbox{\rm \scriptsize law}}{\longrightarrow}}


\renewcommand{\Re}{{\mathfrak{Re}}}
\renewcommand{\Im}{{\mathfrak{Im}}}

% THEOREMS -------------------------------------------------------

\newtheorem{thm}{Theorem}[section]
\newtheorem{cor}[thm]{Corollary}
\newtheorem{lem}[thm]{Lemma}
\newtheorem{prop}[thm]{Proposition}
\theoremstyle{definition}
\newtheorem{defn}[thm]{Definition}
\theoremstyle{remark}
\newtheorem*{rem}{Remark}



\numberwithin{equation}{section}



\begin{document}
\begin{center}
{\bf Probability, homework 5, due October 14.}
\end{center}
\vspace{0.5cm}
\small

\noindent{\bf Exercise 1}.
Prove that if a sequence of real random variables $(X_n)$ converge in distribution to $X$, and $(Y_n)$ converges in distribution to a constant $c$, then $X_n+Y_n$ converges in distribution to $X+c$.\\


\noindent{\bf Exercise 2.}
Assume that $(X,Y)$ has joint density 
$$
c e^{-(1+x^2)(1+y^2)},
$$
where $c$ is properly chosen.
Are $X$ and $Y$ Gaussian random variables? Is $(X,Y)$ a Gaussian vector?
\\




\noindent{\bf Exercise 3}. You toss a coin repeatedly and independently. The probability to get a head is $p$, a tail is $1-p$.
Let $A_k$ be the following event: $k$ or more consecutive heads occur amongst the tosses numbered 
$2^k,\dots,2^{k+1}-1$. Prove that
$\Prob(A_k\ {\rm i.o.})=1$ if $p\geq 1/2$, $0$ otherwise.

Here, i.o. stands for ``infinitely often'', and $A_k$ i.o. is the event $\cap_{n\geq 1} \cup_{m\geq n} A_m$.\\


\noindent{\bf Exercise 4}.  Prove that for any $x>0$, $\frac{1}{x}=\int e^{-tx}\rd t$.  Deduce the value of $\int_0^{\infty}\frac{\sin x}{x}\rd x$.\\





\noindent{\bf Exercise 5}.  For any probability measure $\mu$ supported on $\mathbb{R}_+$,  one defines the Laplace transform as
$$
\mathscr{L}_\mu(\lambda)=\int_0^\infty e^{-\lambda x} \rd \mu(x),\ \lambda\geq 0.
$$
\begin{enumerate}[(1)]
\item Prove that $\mathscr{L}_\mu$ is well-defined,  continuous on $\mathbb{R}_+$ and $\mathscr{C}^\infty$ on 
$\mathbb{R}_+^*$.
\item Prove that $\mathscr{L}_\mu$ characterizes the probability measure $\mu$  supported on $\mathbb{R}_+$.
\item Assume that for a sequence $(\mu_n)_{n\geq 1}$ of probability measure  supported on $\mathbb{R}_+$, one has $\mathscr{L}_{\mu_n}(\lambda)\to \ell(\lambda)$ for any $\lambda\geq 0$, and $\ell$ is right-continuous at $0$.
Prove that $(\mu_n)_{n\geq 1}$ is tight,  and that it converges weakly to a measure $\mu$ such that $\ell=\mathscr{L}_\mu$.
\end{enumerate}




\end{document} 